%===============================================================
% Notas del curso de Algebra Lineal, unidad 3. Primera propuesta.
% Tecnologico Nacional de Mexico, ACF-0903, agosto-diciembre 2026.
% Archivo autonomo: compilar con pdflatex (dos veces) o latexmk.
%===============================================================
\documentclass[12pt]{article}
\usepackage[utf8]{inputenc}
\usepackage[T1]{fontenc}
\usepackage[spanish,mexico]{babel}
\usepackage{lmodern}
\usepackage{amsmath,amssymb,amsthm,mathtools}
\usepackage{geometry,enumitem,booktabs,array,xcolor,graphicx,hyperref,fancyhdr,microtype,longtable,multicol}
\geometry{letterpaper,top=2.4cm,bottom=2.35cm,left=2.45cm,right=2.45cm}
\setlength{\parindent}{0pt}
\setlength{\parskip}{0.52em}
\setlength{\headheight}{15pt}
\setlength{\emergencystretch}{2em}
\setlist[itemize]{leftmargin=1.6em}
\setlist[enumerate]{leftmargin=1.9em}
\definecolor{azul}{RGB}{22,72,120}
\hypersetup{colorlinks=true,linkcolor=azul,urlcolor=azul,pdftitle={Algebra Lineal - Unidad 3: Matrices y determinantes},pdfauthor={Carlos Ernesto Martinez}}
\pagestyle{fancy}
\fancyhf{}
\lhead{\small Tecnológico Nacional de México}
\rhead{\small Álgebra Lineal}
\lfoot{\footnotesize Carlos Ernesto Martínez}
\cfoot{\footnotesize\thepage}
\rfoot{\footnotesize mrdgzcarlose@gmail.com}
\renewcommand{\headrulewidth}{0.4pt}
\renewcommand{\footrulewidth}{0.3pt}
\newtheorem{definicion}{Definición}[section]
\newtheorem{teorema}[definicion]{Teorema}
\newtheorem{proposicion}[definicion]{Proposición}
\newtheorem{corolario}[definicion]{Corolario}
\theoremstyle{remark}
\newtheorem{observacion}[definicion]{Observación}
\newtheorem{ejemplo}[definicion]{Ejemplo}
\newcommand{\R}{\mathbb{R}}
\newcommand{\C}{\mathbb{C}}
\newcommand{\Mat}[2]{\mathcal{M}_{#1\times #2}}
\newcommand{\rang}{\operatorname{rang}}
\newcommand{\nul}{\operatorname{nul}}
\newcommand{\tr}{\operatorname{tr}}
\newcommand{\adj}{\operatorname{adj}}
\newcommand{\Ker}{\operatorname{Ker}}
\newcommand{\RREF}{\operatorname{RREF}}
\newcommand{\Span}{\operatorname{span}}
\newcommand{\vect}[1]{\boldsymbol{#1}}
\newcommand{\idea}[1]{\begin{center}\fcolorbox{azul}{white}{\parbox{0.91\linewidth}{#1}}\end{center}}
\newcounter{contadorEjercicio}
\newenvironment{listaejercicios}{\begin{enumerate}[label=\textbf{E\arabic*.},leftmargin=2.4em,itemsep=0.52em]\setcounter{enumi}{\value{contadorEjercicio}}}{\setcounter{contadorEjercicio}{\value{enumi}}\end{enumerate}}
\title{\textbf{Notas para el curso de Álgebra Lineal}\\[0.35cm]
\Large Unidad 3: Matrices y determinantes\\[0.25cm]
\normalsize Primera propuesta de documento de trabajo}
\author{Carlos Ernesto Martínez\\Tecnológico Nacional de México\\Clave ACF-0903}
\date{Agosto--Diciembre 2026}
\begin{document}
\maketitle
\begin{center}\textcolor{azul}{\rule{0.72\textwidth}{0.6pt}}\end{center}
\begin{abstract}
Estas notas desarrollan la Unidad 3 del temario de Álgebra Lineal (ACF-0903): concepto y tipos de matrices, operaciones, transposición, inversas, determinantes, menores y cofactores, adjunta, rango y núcleo. Se retoman las operaciones elementales y el método de Gauss--Jordan estudiados en la Unidad 2, sin repetir toda su teoría. Las demostraciones, los ejemplos resueltos y los ejercicios están pensados para uso en clase y estudio independiente.
\end{abstract}
\tableofcontents
\newpage

%===============================================================
\section{Propósito y organización de la unidad}
%===============================================================
En la Unidad 2 utilizamos matrices para representar sistemas de ecuaciones lineales y realizar operaciones elementales por renglones. Ahora estudiaremos las matrices como objetos matemáticos por derecho propio. Nos interesa comprender cuáles son sus operaciones, qué propiedades satisfacen, cuándo es posible invertirlas y cómo los determinantes aportan información sobre una matriz cuadrada.

Finalmente, introduciremos el rango y el núcleo, conceptos que aparecerán de nuevo en las Unidades 4 y 5, al estudiar espacios vectoriales y transformaciones lineales.

\idea{\textbf{Ruta conceptual:}\quad matriz y sus tipos $\longrightarrow$ operaciones $\longrightarrow$ inversa $\longrightarrow$ determinante y adjunta $\longrightarrow$ rango y núcleo.}

\textbf{Contenidos del temario.} Esta unidad cubre los 27 subtemas de la Unidad 3 del documento de trabajo del curso: definición y clasificación de matrices, operaciones, transposición, inversa, determinantes, menores, cofactores, adjunta, rango, núcleo y aplicaciones matriciales.

\textbf{Tiempo previsto:} 10 horas (6 teóricas y 4 prácticas), distribuidas en dos semanas. Los ejercicios y demostraciones adicionales permiten adaptar el ritmo de trabajo y la profundidad según el grupo.

\subsection*{Conocimientos previos}
Se supondrá que el estudiante conoce las operaciones algebraicas usuales, los sistemas de ecuaciones lineales y las operaciones elementales por renglones. De forma predeterminada, trabajaremos con matrices reales; cuando corresponda, indicaremos qué afirmaciones también son válidas para matrices complejas.

%===============================================================
\section{Matrices: definición, orden y clasificación}
%===============================================================
\subsection{Definición y notación}
Una matriz permite organizar datos o coeficientes mediante filas y columnas. El objeto resultante no es sólo una tabla: su estructura hace posible definir operaciones algebraicas con propiedades precisas.

\begin{definicion}[Matriz]
Una \emph{matriz real de orden $m\times n$} es un arreglo rectangular de $mn$ números reales, distribuidos en $m$ filas y $n$ columnas:
\[
A=\begin{pmatrix}
a_{11}&a_{12}&\cdots&a_{1n}\\
a_{21}&a_{22}&\cdots&a_{2n}\\
\vdots&\vdots&\ddots&\vdots\\
a_{m1}&a_{m2}&\cdots&a_{mn}
\end{pmatrix}=\bigl(a_{ij}\bigr)_{m\times n}.
\]
El número $a_{ij}$ es la entrada situada en la fila $i$ y la columna $j$. El conjunto de todas las matrices reales $m\times n$ se denota por $\Mat{m}{n}(\R)$.
\end{definicion}

\begin{ejemplo}
Sea
\[
A=\begin{pmatrix}2&-1&3\\0&4&5\end{pmatrix}.
\]
La matriz $A$ tiene dos filas y tres columnas, así que su orden es $2\times3$. Las entradas $a_{12}=-1$ y $a_{23}=5$ indican, respectivamente, el elemento de la primera fila y segunda columna, y el de la segunda fila y tercera columna.
\end{ejemplo}

\begin{definicion}[Igualdad de matrices]
Dos matrices $A=(a_{ij})$ y $B=(b_{ij})$ son iguales si y sólo si tienen el mismo orden y sus entradas correspondientes coinciden:
\[
A=B\quad\Longleftrightarrow\quad a_{ij}=b_{ij}\quad\text{para todo }i,j.
\]
\end{definicion}

\begin{ejemplo}
Encontrar $x,y,z$ si
\[
\begin{pmatrix}x+1&2y\\z-3&5\end{pmatrix}=
\begin{pmatrix}4&-6\\1&5\end{pmatrix}.
\]
Igualamos las entradas correspondientes:
\[
x+1=4,\qquad 2y=-6,\qquad z-3=1.
\]
Por tanto, $\boxed{x=3,\ y=-3,\ z=4}$.
\end{ejemplo}

\subsection{Tipos de matrices}
\begin{definicion}[Matrices especiales]
Consideremos una matriz $A=(a_{ij})$.
\begin{itemize}
\item \textbf{Fila:} tiene una sola fila, es decir, orden $1\times n$.
\item \textbf{Columna:} tiene una sola columna, orden $m\times1$.
\item \textbf{Rectangular:} su número de filas es distinto del de columnas ($m\ne n$).
\item \textbf{Cuadrada:} tiene el mismo número de filas y columnas ($n\times n$).
\item \textbf{Nula:} todas sus entradas son cero; se denota por $0$ cuando el orden está claro.
\item \textbf{Diagonal:} es cuadrada y $a_{ij}=0$ cuando $i\ne j$.
\item \textbf{Identidad:} es diagonal, con $a_{ii}=1$ para todo $i$; se denota por $I_n$.
\item \textbf{Triangular superior:} $a_{ij}=0$ si $i>j$.
\item \textbf{Triangular inferior:} $a_{ij}=0$ si $i<j$.
\item \textbf{Simétrica:} $A=A^T$.
\end{itemize}
\end{definicion}

\begin{ejemplo}
Clasifiquemos las siguientes matrices:
\[
A=\begin{pmatrix}2&0&0\\0&-1&0\\0&0&4\end{pmatrix},\qquad
B=\begin{pmatrix}1&3&-2\\0&5&4\\0&0&7\end{pmatrix},\qquad
C=\begin{pmatrix}2&-1\\-1&3\end{pmatrix}.
\]
$A$ es diagonal (y, simultáneamente, triangular superior e inferior); $B$ es triangular superior, pero no diagonal; $C$ es simétrica porque sus entradas fuera de la diagonal coinciden al intercambiar filas y columnas.
\end{ejemplo}

\begin{observacion}
Una matriz diagonal es siempre simétrica. La recíproca es falsa: por ejemplo, $\begin{psmallmatrix}1&2\\2&3\end{psmallmatrix}$ es simétrica y no es diagonal. Además, una matriz puede pertenecer a varias categorías a la vez.
\end{observacion}

%===============================================================
\section{Operaciones con matrices}
%===============================================================
\subsection{Suma y multiplicación por un escalar}
\begin{definicion}[Suma]
Para matrices $A=(a_{ij})$ y $B=(b_{ij})$ del \emph{mismo orden}, definimos
\[
A+B=(a_{ij}+b_{ij}).
\]
La suma no está definida, bajo esta operación, si los órdenes de $A$ y $B$ son diferentes.
\end{definicion}

\begin{definicion}[Producto por un escalar]
Sea $\lambda\in\R$ y $A=(a_{ij})$. Definimos
\[
\lambda A=(\lambda a_{ij}).
\]
El orden de $\lambda A$ es el mismo de $A$.
\end{definicion}

\begin{ejemplo}
Sean $A=\begin{psmallmatrix}1&-2\\3&0\end{psmallmatrix}$ y $B=\begin{psmallmatrix}4&1\\-1&5\end{psmallmatrix}$. Entonces
\[
A+B=\begin{pmatrix}5&-1\\2&5\end{pmatrix},\qquad
2A-B=\begin{pmatrix}-2&-5\\7&-5\end{pmatrix}.
\]
En $2A-B$ multiplicamos primero cada entrada de $A$ por $2$ y después restamos las entradas correspondientes de $B$.
\end{ejemplo}

\begin{proposicion}[Propiedades aditivas]
Para $A,B,C\in\Mat{m}{n}(\R)$ y $\lambda,\mu\in\R$ se cumple
\begin{align*}
A+B&=B+A,& (A+B)+C&=A+(B+C),\\
A+0&=A,& A+(-A)&=0,\\
\lambda(A+B)&=\lambda A+\lambda B,& (\lambda+\mu)A&=\lambda A+\mu A,\\
(\lambda\mu)A&=\lambda(\mu A),&1A&=A.
\end{align*}
\end{proposicion}
\begin{proof}
La demostración es \emph{entrada por entrada}. Por ejemplo,
\[
\bigl(\lambda(A+B)\bigr)_{ij}=\lambda(a_{ij}+b_{ij})
=\lambda a_{ij}+\lambda b_{ij}
=\bigl(\lambda A+\lambda B\bigr)_{ij}.
\]
Las demás identidades se obtienen aplicando las propiedades usuales de los números reales a cada entrada.
\end{proof}

\subsection{Producto de matrices}
El producto matricial no se calcula multiplicando simplemente entradas que ocupen la misma posición. La regla combina las filas de la primera matriz con las columnas de la segunda.

\begin{definicion}[Producto matricial]
Sean $A=(a_{ij})\in\Mat{m}{n}(\R)$ y $B=(b_{jk})\in\Mat{n}{p}(\R)$. Se define $C=AB\in\Mat{m}{p}(\R)$ mediante
\[
\boxed{\quad c_{ik}=\sum_{j=1}^{n}a_{ij}b_{jk}\quad}
\]
para $1\le i\le m$, $1\le k\le p$.
\end{definicion}

La condición de existencia del producto es que el número de columnas de $A$ coincida con el número de filas de $B$:
\[
\underbrace{(m\times n)}_A\ \underbrace{(n\times p)}_B
\longrightarrow\underbrace{(m\times p)}_{AB}.
\]

\begin{ejemplo}[Producto de una matriz $2\times3$ por una $3\times2$]
Calculemos $AB$ para
\[
A=\begin{pmatrix}1&2&-1\\0&3&4\end{pmatrix},\qquad
B=\begin{pmatrix}2&1\\-1&0\\3&2\end{pmatrix}.
\]
El producto existe y tiene orden $2\times2$. Calculamos cada entrada:
\begin{align*}
c_{11}&=1(2)+2(-1)+(-1)(3)=-3,\\
c_{12}&=1(1)+2(0)+(-1)(2)=-1,\\
c_{21}&=0(2)+3(-1)+4(3)=9,\\
c_{22}&=0(1)+3(0)+4(2)=8.
\end{align*}
Por tanto,
\[\boxed{AB=\begin{pmatrix}-3&-1\\9&8\end{pmatrix}.}\]
\end{ejemplo}

\begin{observacion}
Aunque el producto $AB$ exista, $BA$ podría no existir. Cuando ambos existen, generalmente tienen órdenes distintos o no son iguales.
\end{observacion}

\begin{ejemplo}[El producto no es conmutativo]
Sean
\[
A=\begin{pmatrix}1&1\\0&1\end{pmatrix},\qquad
B=\begin{pmatrix}1&0\\1&1\end{pmatrix}.
\]
Entonces
\[
AB=\begin{pmatrix}2&1\\1&1\end{pmatrix},\qquad
BA=\begin{pmatrix}1&1\\1&2\end{pmatrix}.
\]
Ambos productos existen, pero $AB\ne BA$. En consecuencia, \emph{no} podemos usar indiscriminadamente las reglas conmutativas de los números reales.
\end{ejemplo}

\begin{proposicion}[Propiedades del producto]
Si las dimensiones permiten realizar las operaciones, entonces
\begin{align*}
A(BC)&=(AB)C,\\
A(B+C)&=AB+AC,\\
(A+B)C&=AC+BC,\\
(\lambda A)B&=\lambda(AB)=A(\lambda B),\\
I_mA&=A=AI_n,\qquad A\in\Mat{m}{n}(\R).
\end{align*}
\end{proposicion}
\begin{proof}[Demostración de la asociatividad]
Sean $A$ de orden $m\times n$, $B$ de orden $n\times p$ y $C$ de orden $p\times q$. La entrada $(i,k)$ de $(AB)C$ es
\[
\sum_{\ell=1}^{p}\left(\sum_{j=1}^{n}a_{ij}b_{j\ell}\right)c_{\ell k}
=\sum_{j=1}^{n}\sum_{\ell=1}^{p}a_{ij}b_{j\ell}c_{\ell k}.
\]
La entrada $(i,k)$ de $A(BC)$ es
\[
\sum_{j=1}^{n}a_{ij}\left(\sum_{\ell=1}^{p}b_{j\ell}c_{\ell k}\right)
=\sum_{j=1}^{n}\sum_{\ell=1}^{p}a_{ij}b_{j\ell}c_{\ell k}.
\]
Son iguales porque las sumas son finitas; por tanto, $(AB)C=A(BC)$.
\end{proof}

\subsection{Transpuesta de una matriz}
\begin{definicion}[Matriz transpuesta]
Para $A=(a_{ij})\in\Mat{m}{n}(\R)$, su \emph{transpuesta} es la matriz $A^T\in\Mat{n}{m}(\R)$ definida por
\[
(A^T)_{ji}=a_{ij}.
\]
En otras palabras, las filas se convierten en columnas.
\end{definicion}
\begin{ejemplo}
Si $A=\begin{psmallmatrix}1&-2&3\\4&0&5\end{psmallmatrix}$, entonces
\[
A^T=\begin{pmatrix}1&4\\-2&0\\3&5\end{pmatrix}.
\]
\end{ejemplo}

\begin{proposicion}[Propiedades de la transposición]
Para matrices de tamaños compatibles y $\lambda\in\R$:
\[
(A^T)^T=A,\qquad(A+B)^T=A^T+B^T,\qquad
(\lambda A)^T=\lambda A^T,\qquad (AB)^T=B^TA^T.
\]
\end{proposicion}
\begin{proof}[Demostración de $(AB)^T=B^TA^T$]
Sean $A\in\Mat{m}{n}(\R)$ y $B\in\Mat{n}{p}(\R)$. Para cada $i,k$:
\[
\bigl((AB)^T\bigr)_{ki}=(AB)_{ik}
=\sum_{j=1}^{n}a_{ij}b_{jk}
=\sum_{j=1}^{n}(B^T)_{kj}(A^T)_{ji}
=(B^TA^T)_{ki}.
\]
Por igualdad de entradas, la identidad queda demostrada.
\end{proof}

\begin{observacion}
Para matrices con entradas complejas, además de la transpuesta se utiliza la \emph{transpuesta conjugada} $A^*=\overline{A}^{\,T}$. La distinción será útil posteriormente al hablar de producto interno. No necesitamos introducirla aún para desarrollar los contenidos obligatorios de esta unidad.
\end{observacion}

%===============================================================
\section{Matrices inversas y el método de Gauss--Jordan}
%===============================================================
\subsection{Concepto de inversa}
\begin{definicion}[Matriz invertible]
Una matriz cuadrada $A\in\Mat{n}{n}(\R)$ es \emph{invertible} si existe una matriz $B$ del mismo orden tal que
\[
AB=BA=I_n.
\]
La matriz $B$ se llama \emph{inversa} de $A$ y se escribe $A^{-1}$. Si no existe, $A$ es \emph{singular} o \emph{no invertible}.
\end{definicion}

\begin{proposicion}[Unicidad de la inversa]
Si $A$ es invertible, su inversa es única.
\end{proposicion}
\begin{proof}
Supongamos que $B$ y $C$ son inversas de $A$. Usando asociatividad,
\[
B=BI_n=B(AC)=(BA)C=I_nC=C.
\]
Por tanto, no pueden existir dos inversas distintas.
\end{proof}

\begin{proposicion}[Inversa de un producto]
Si $A$ y $B$ son matrices cuadradas invertibles del mismo orden, entonces
\[
\boxed{(AB)^{-1}=B^{-1}A^{-1}.}
\]
\end{proposicion}
\begin{proof}
Basta comprobar
\[
(AB)(B^{-1}A^{-1})=A(BB^{-1})A^{-1}=AA^{-1}=I_n,
\]
y, análogamente, $(B^{-1}A^{-1})(AB)=I_n$. La unicidad de la inversa concluye la demostración.
\end{proof}

\subsection{Obtención de la inversa mediante Gauss--Jordan}
La idea consiste en aplicar operaciones elementales simultáneamente a $A$ y a la identidad, mediante la matriz aumentada
\[
\left[A\mid I_n\right].
\]
Si las operaciones convierten el bloque izquierdo en $I_n$, el bloque derecho resulta ser $A^{-1}$:
\[
\boxed{[A\mid I_n]\ \sim\ [I_n\mid A^{-1}].}
\]
Cuando el bloque izquierdo no puede reducirse a la identidad, $A$ no es invertible.

\begin{ejemplo}[Cálculo completo por Gauss--Jordan]
Sea
\[
A=\begin{pmatrix}2&1\\1&1\end{pmatrix}.
\]
Formamos la matriz aumentada y efectuamos las operaciones indicadas:
\begin{align*}
\left(\begin{array}{cc|cc}2&1&1&0\\1&1&0&1\end{array}\right)
&\xrightarrow{R_1\leftrightarrow R_2}
\left(\begin{array}{cc|cc}1&1&0&1\\2&1&1&0\end{array}\right)\\[1ex]
&\xrightarrow{R_2\leftarrow R_2-2R_1}
\left(\begin{array}{cc|cc}1&1&0&1\\0&-1&1&-2\end{array}\right)\\[1ex]
&\xrightarrow{R_2\leftarrow -R_2}
\left(\begin{array}{cc|cc}1&1&0&1\\0&1&-1&2\end{array}\right)\\[1ex]
&\xrightarrow{R_1\leftarrow R_1-R_2}
\left(\begin{array}{cc|cc}1&0&1&-1\\0&1&-1&2\end{array}\right).
\end{align*}
Leemos la matriz inversa del bloque derecho:
\[
\boxed{A^{-1}=\begin{pmatrix}1&-1\\-1&2\end{pmatrix}.}
\]
Verificamos multiplicando:
\[
AA^{-1}=\begin{pmatrix}2&1\\1&1\end{pmatrix}
\begin{pmatrix}1&-1\\-1&2\end{pmatrix}
=\begin{pmatrix}1&0\\0&1\end{pmatrix}.
\]
\end{ejemplo}

\begin{ejemplo}[Matriz que no tiene inversa]
Sea
\[
B=\begin{pmatrix}1&2\\2&4\end{pmatrix}.
\]
La segunda fila es el doble de la primera. La operación $R_2\leftarrow R_2-2R_1$ crea una fila nula en el bloque izquierdo. Por ello el bloque izquierdo nunca puede convertirse en $I_2$ mediante operaciones elementales; $B$ no es invertible.
\end{ejemplo}

\begin{observacion}
Resolver $A\vect{x}=\vect{b}$ mediante $\vect{x}=A^{-1}\vect{b}$ exige que $A$ sea cuadrada e invertible. En problemas grandes suele ser preferible resolver directamente el sistema sin construir explícitamente $A^{-1}$.
\end{observacion}

%===============================================================
\section{Determinantes}
%===============================================================
\subsection{Determinantes de orden dos y tres}
El determinante asocia a cada matriz cuadrada un número. Este número permite distinguir, entre otras cosas, matrices invertibles de matrices singulares.

\begin{definicion}[Determinante de orden dos]
Para $A=\begin{psmallmatrix}a&b\\c&d\end{psmallmatrix}$ definimos
\[
\det(A)=\begin{vmatrix}a&b\\c&d\end{vmatrix}=ad-bc.
\]
\end{definicion}
\begin{ejemplo}
Para $A=\begin{psmallmatrix}2&1\\1&1\end{psmallmatrix}$ tenemos
\[
\det(A)=2(1)-1(1)=1.
\]
El determinante no nulo anticipa que $A$ es invertible, como ya comprobamos por Gauss--Jordan.
\end{ejemplo}

\begin{definicion}[Determinante de orden tres]
Si
\[
A=\begin{pmatrix}a&b&c\\d&e&f\\g&h&i\end{pmatrix},
\]
su determinante es
\[
\boxed{\det(A)=aei+bfg+cdh-ceg-bdi-afh.}
\]
Esta expresión puede organizarse mediante la llamada \emph{regla de Sarrus}, válida únicamente para matrices $3\times3$.
\end{definicion}

\begin{ejemplo}
Calcular
\[
D=\begin{vmatrix}1&2&0\\0&3&1\\2&0&1\end{vmatrix}.
\]
Por la fórmula anterior:
\begin{align*}
D&=(1)(3)(1)+(2)(1)(2)+(0)(0)(0)\\
&\quad -(0)(3)(2)-(2)(0)(1)-(1)(1)(0)\\
&=3+4=\boxed{7}.
\end{align*}
\end{ejemplo}

\subsection{Definición general: menores y cofactores}
\begin{definicion}[Menor y cofactor]
Sea $A=(a_{ij})$ una matriz cuadrada de orden $n\ge2$. El \emph{menor} $M_{ij}$ de la entrada $a_{ij}$ es el determinante de la matriz de orden $n-1$ obtenida al suprimir la fila $i$ y la columna $j$.

El \emph{cofactor} correspondiente es
\[
C_{ij}=(-1)^{i+j}M_{ij}.
\]
\end{definicion}

Los signos de los cofactores alternan siguiendo el patrón
\[
\begin{pmatrix}
+&-&+&\cdots\\-&+&-&\cdots\\+&-&+&\cdots\\\vdots&\vdots&\vdots&\ddots
\end{pmatrix}.
\]

\begin{teorema}[Expansión por cofactores o de Laplace]
Sea $A=(a_{ij})$ una matriz cuadrada de orden $n\ge2$. Fijando una fila $i$, su determinante puede calcularse como
\[
\boxed{\det(A)=\sum_{j=1}^{n}a_{ij}C_{ij}.}
\]
También puede expandirse a lo largo de cualquier columna $j$:
\[
\det(A)=\sum_{i=1}^{n}a_{ij}C_{ij}.
\]
Para cerrar la definición recursiva se toma $\det([a])=a$ para matrices de orden uno.
\end{teorema}

\begin{observacion}
Para calcular un determinante mediante cofactores conviene elegir una fila o una columna con muchos ceros: se reduce así el número de determinantes menores que hay que evaluar.
\end{observacion}

\begin{ejemplo}[Expansión por la primera fila]
Consideremos nuevamente
\[
A=\begin{pmatrix}1&2&0\\0&3&1\\2&0&1\end{pmatrix}.
\]
Los menores de la primera fila son
\[
M_{11}=\begin{vmatrix}3&1\\0&1\end{vmatrix}=3,\quad
M_{12}=\begin{vmatrix}0&1\\2&1\end{vmatrix}=-2,\quad
M_{13}=\begin{vmatrix}0&3\\2&0\end{vmatrix}=-6.
\]
Los cofactores son $C_{11}=3$, $C_{12}=2$ y $C_{13}=-6$. Por tanto,
\[
\det(A)=1(3)+2(2)+0(-6)=\boxed{7},
\]
en concordancia con el resultado obtenido por Sarrus.
\end{ejemplo}

\subsection{Propiedades de los determinantes}
\begin{proposicion}[Propiedades esenciales]
Para matrices cuadradas de orden $n$:
\begin{enumerate}
\item $\det(I_n)=1$.
\item $\det(A^T)=\det(A)$.
\item Si $A$ es triangular, $\det(A)=\prod_{i=1}^{n}a_{ii}$.
\item Si dos filas de $A$ son iguales o proporcionales, entonces $\det(A)=0$.
\item Intercambiar dos filas cambia el signo del determinante.
\item Multiplicar una fila por $\lambda$ multiplica el determinante por $\lambda$.
\item Sumar a una fila un múltiplo de otra fila \emph{no cambia} el determinante.
\item $\det(AB)=\det(A)\det(B)$.
\end{enumerate}
\end{proposicion}

\begin{proof}[Justificación de la regla de suma de filas]
Una propiedad estructural del determinante es su \emph{linealidad en cada fila}, manteniendo fijas las otras. Si reemplazamos la fila $r_i$ por $r_i+\lambda r_j$, con $i\ne j$, obtenemos la suma de dos determinantes: el original más $\lambda$ veces otro en el que las filas $i$ y $j$ coinciden. El segundo vale cero; por tanto, el determinante no cambia.
\end{proof}

\begin{observacion}
Las propiedades anteriores pueden demostrarse rigurosamente a partir de la definición del determinante mediante permutaciones. Aquí las empleamos como resultados del curso; su justificación completa mediante esa construcción combinatoria sería una ampliación teórica adicional.
\end{observacion}

\begin{ejemplo}[Determinante por eliminación]
Calcular
\[
\det\begin{pmatrix}2&1&0\\4&3&1\\0&2&3\end{pmatrix}.
\]
Aplicamos operaciones que no cambian el determinante:
\[
\begin{pmatrix}2&1&0\\4&3&1\\0&2&3\end{pmatrix}
\xrightarrow{R_2\leftarrow R_2-2R_1}
\begin{pmatrix}2&1&0\\0&1&1\\0&2&3\end{pmatrix}
\xrightarrow{R_3\leftarrow R_3-2R_2}
\begin{pmatrix}2&1&0\\0&1&1\\0&0&1\end{pmatrix}.
\]
La última matriz es triangular superior. Su determinante es el producto de las entradas diagonales:
\[
\boxed{\det(A)=2\cdot1\cdot1=2}.
\]
\end{ejemplo}

\subsection{Determinantes e invertibilidad}
\begin{teorema}[Criterio de invertibilidad]
Para una matriz cuadrada $A$ de orden $n$, las siguientes condiciones son equivalentes:
\begin{enumerate}
\item $A$ es invertible.
\item $\det(A)\ne0$.
\item Las filas de $A$ son linealmente independientes.
\item $\rang(A)=n$.
\item El sistema homogéneo $A\vect{x}=\vect{0}$ sólo tiene la solución trivial.
\end{enumerate}
\end{teorema}
\begin{observacion}
Algunos conceptos de la lista, como independencia lineal, serán estudiados formalmente en la Unidad 4. Por ahora es suficiente relacionar los criterios operativos: pivotes, rango, determinante y existencia de inversa.
\end{observacion}

%===============================================================
\section{Matriz adjunta e inversa mediante cofactores}
%===============================================================
\begin{definicion}[Matriz de cofactores y adjunta]
Sea $A=(a_{ij})$ cuadrada. La \emph{matriz de cofactores} es $C=(C_{ij})$. Su transpuesta recibe el nombre de \emph{matriz adjunta}:
\[
\boxed{\adj(A)=C^T.}
\]
\end{definicion}

\begin{teorema}[Identidad fundamental de la adjunta]
Para cualquier matriz cuadrada $A$,
\[
A\adj(A)=\adj(A)A=\det(A)I_n.
\]
En consecuencia, si $\det(A)\ne0$,
\[
\boxed{A^{-1}=\frac{1}{\det(A)}\adj(A).}
\]
\end{teorema}
\begin{proof}
La entrada $(i,k)$ de $A\adj(A)$ es
\[
\sum_{j=1}^{n}a_{ij}C_{kj}.
\]
Si $i=k$, esa suma es la expansión por cofactores de la fila $i$ y vale $\det(A)$. Si $i\ne k$, la misma suma corresponde al determinante de una matriz en la que se reemplaza la fila $k$ por la fila $i$. Esa matriz tiene dos filas iguales y su determinante es cero. Por ello $A\adj(A)=\det(A)I_n$. La identidad por el otro lado se obtiene razonando con columnas. Si el determinante no es cero, dividimos entre $\det(A)$ para obtener la inversa.
\end{proof}

\begin{ejemplo}[Inversa $2\times2$ mediante la adjunta]
Para
\[
A=\begin{pmatrix}a&b\\c&d\end{pmatrix},\qquad \det(A)=ad-bc\ne0,
\]
la matriz de cofactores es $\begin{psmallmatrix}d&-c\\-b&a\end{psmallmatrix}$ y su adjunta es $\begin{psmallmatrix}d&-b\\-c&a\end{psmallmatrix}$. Así,
\[
\boxed{A^{-1}=\frac{1}{ad-bc}\begin{pmatrix}d&-b\\-c&a\end{pmatrix}.}
\]
En particular, para $A=\begin{psmallmatrix}2&1\\1&1\end{psmallmatrix}$:
\[
\det(A)=1,\qquad A^{-1}=\begin{pmatrix}1&-1\\-1&2\end{pmatrix},
\]
el mismo resultado obtenido por Gauss--Jordan.
\end{ejemplo}

\begin{ejemplo}[Adjunta e inversa de una matriz $3\times3$]
Consideremos
\[
A=\begin{pmatrix}2&1&0\\1&2&1\\0&1&2\end{pmatrix}.
\]
Expandiendo por la primera fila:
\[
\det(A)=2\begin{vmatrix}2&1\\1&2\end{vmatrix}
-1\begin{vmatrix}1&1\\0&2\end{vmatrix}
=2(3)-2=4.
\]
Calculamos los nueve cofactores:
\[
C=\begin{pmatrix}
3&-2&1\\
-2&4&-2\\
1&-2&3
\end{pmatrix}.
\]
En este ejemplo $C=C^T$, por lo que
\[
\adj(A)=C^T=\begin{pmatrix}
3&-2&1\\-2&4&-2\\1&-2&3
\end{pmatrix}.
\]
Se concluye
\[
\boxed{A^{-1}=\frac14\begin{pmatrix}
3&-2&1\\-2&4&-2\\1&-2&3
\end{pmatrix}.}
\]
Comprobación, multiplicando la primera fila de $A$ por las columnas de la adjunta:
\[
(2,1,0)\adj(A)=(4,0,0).
\]
Por simetría y multiplicación directa en las otras filas, $A\adj(A)=4I_3$, como afirma el teorema.
\end{ejemplo}

%===============================================================
\section{Rango y núcleo de una matriz}
%===============================================================
Estos conceptos sirven como puente hacia los espacios vectoriales y las transformaciones lineales. En esta unidad los introduciremos a partir del cálculo con matrices.

\begin{definicion}[Rango]
El \emph{rango} de una matriz $A$, denotado por $\rang(A)$, es el número de pivotes de una forma escalonada obtenida mediante operaciones elementales por filas. Equivalentemente, es el número de columnas linealmente independientes de $A$.
\end{definicion}

\begin{definicion}[Núcleo]
Para $A\in\Mat{m}{n}(\R)$, su \emph{núcleo} es el conjunto de soluciones del sistema homogéneo asociado:
\[
\boxed{\Ker(A)=\{\vect{x}\in\R^n:A\vect{x}=\vect{0}\}.}
\]
\end{definicion}

\begin{observacion}
Las operaciones elementales por filas preservan las soluciones de $A\vect{x}=\vect0$. Por tanto, podemos calcular el núcleo a partir de la forma escalonada reducida. Sin embargo, esas operaciones no preservan, en general, el \emph{espacio de columnas} como subconjunto de $\R^m$; para construir una base de las columnas utilizaremos las columnas pivote de la \emph{matriz original}.
\end{observacion}

\begin{ejemplo}[Rango y núcleo paso a paso]
Sea
\[
A=\begin{pmatrix}1&2&3\\2&4&6\\0&1&1\end{pmatrix}.
\]
Buscamos la forma escalonada reducida:
\begin{align*}
A&\xrightarrow{R_2\leftarrow R_2-2R_1}
\begin{pmatrix}1&2&3\\0&0&0\\0&1&1\end{pmatrix}
\xrightarrow{R_2\leftrightarrow R_3}
\begin{pmatrix}1&2&3\\0&1&1\\0&0&0\end{pmatrix}\\[1ex]
&\xrightarrow{R_1\leftarrow R_1-2R_2}
\begin{pmatrix}1&0&1\\0&1&1\\0&0&0\end{pmatrix}.
\end{align*}
Hay dos pivotes; por tanto,
\[\boxed{\rang(A)=2.}\]
El sistema $A\vect{x}=\vect{0}$ equivale a
\[
\begin{cases}
x_1+x_3=0,\\x_2+x_3=0.
\end{cases}
\]
Tomando $x_3=t$ como variable libre, obtenemos
\[
x_1=-t,\qquad x_2=-t,\qquad x_3=t.
\]
Luego
\[
\boxed{\Ker(A)=\left\{t\begin{pmatrix}-1\\-1\\1\end{pmatrix}:t\in\R\right\}.}
\]
Las columnas pivote son la primera y la segunda; en la matriz original son
\[
\begin{pmatrix}1\\2\\0\end{pmatrix},\qquad
\begin{pmatrix}2\\4\\1\end{pmatrix}.
\]
La tercera es la suma de las dos primeras. Esta observación da una interpretación concreta de $\rang(A)=2$.
\end{ejemplo}

\begin{proposicion}
Para toda matriz $A\in\Mat{m}{n}(\R)$,
\[
0\le\rang(A)\le\min\{m,n\}.
\]
Si $A$ es cuadrada de orden $n$, entonces es invertible si y sólo si $\rang(A)=n$.
\end{proposicion}

\begin{observacion}[Anticipo de la Unidad 5]
Si el núcleo de una matriz $m\times n$ tiene dimensión $d$, el teorema rango--nulidad establece que $\rang(A)+d=n$. En la Unidad 5 demostraremos este resultado en el lenguaje de las transformaciones lineales.
\end{observacion}

%===============================================================
\section{Aplicación elemental del álgebra matricial}
%===============================================================
\begin{ejemplo}[Modelo sencillo mediante un sistema matricial]
Dos cantidades $x$ e $y$ satisfacen
\[
\begin{cases}2x+y=5,\\x+y=3.\end{cases}
\]
En forma matricial,
\[
\underbrace{\begin{pmatrix}2&1\\1&1\end{pmatrix}}_{A}
\underbrace{\begin{pmatrix}x\\y\end{pmatrix}}_{\vect{x}}
=\underbrace{\begin{pmatrix}5\\3\end{pmatrix}}_{\vect{b}}.
\]
Como $\det(A)=1$, $A$ es invertible. Usando la inversa calculada anteriormente:
\[
\vect{x}=A^{-1}\vect{b}=
\begin{pmatrix}1&-1\\-1&2\end{pmatrix}
\begin{pmatrix}5\\3\end{pmatrix}
=\begin{pmatrix}2\\1\end{pmatrix}.
\]
Por tanto, $\boxed{x=2,\ y=1}$. Sustituyendo en las dos ecuaciones originales verificamos que ambas se satisfacen.
\end{ejemplo}

\begin{observacion}
Esta aplicación algebraica se presenta deliberadamente sin un contexto artificial. Para otras aplicaciones en ingeniería se requerirá formular variables, unidades y supuestos del modelo; el procedimiento matricial empieza sólo después de ese planteamiento.
\end{observacion}

%===============================================================
\section{Banco de ejercicios}
%===============================================================
Los ejercicios están ordenados por tema. Se recomienda resolverlos primero sin software y, después, utilizar R, Python o una calculadora como instrumento de verificación. El apéndice contiene soluciones de una selección representativa.

\subsection{Definición, clasificación e igualdad}
\begin{listaejercicios}
\item Determine el orden de $A=\begin{psmallmatrix}1&0&2\\-3&4&5\end{psmallmatrix}$ e identifique $a_{12}$, $a_{23}$.
\item Clasifique $A=\begin{psmallmatrix}3&0&0\\0&3&0\\0&0&3\end{psmallmatrix}$; indique todas las categorías de la Unidad 3 a las que pertenece.
\item Determine $x,y,z$ si $\begin{psmallmatrix}x-2&3y\\z+1&7\end{psmallmatrix}=\begin{psmallmatrix}5&-6\\4&7\end{psmallmatrix}$.
\item Dé un ejemplo de matriz simétrica $3\times3$ que no sea diagonal y de matriz triangular inferior no diagonal.
\item Decida cuáles de las matrices $\begin{psmallmatrix}1&2\\2&1\end{psmallmatrix}$, $\begin{psmallmatrix}1&0\\3&1\end{psmallmatrix}$ e $I_2$ son simétricas.
\end{listaejercicios}

\subsection{Operaciones y transposición}
\begin{listaejercicios}
\item Sean $A=\begin{psmallmatrix}1&2\\3&4\end{psmallmatrix}$ y $B=\begin{psmallmatrix}-1&0\\2&5\end{psmallmatrix}$. Calcule $A+B$, $2A-B$ y $A^T$.
\item Calcule $AB$ y $BA$ para $A=\begin{psmallmatrix}1&2\\0&1\end{psmallmatrix}$ y $B=\begin{psmallmatrix}2&0\\1&3\end{psmallmatrix}$.
\item Si $A$ es $2\times3$, $B$ es $3\times4$ y $C$ es $4\times2$, indique el orden de $AB$, $BC$, $(AB)C$ y $C(AB)$.
\item Calcule $AB$ para $A=\begin{psmallmatrix}2&-1&0\\1&3&2\end{psmallmatrix}$ y $B=\begin{psmallmatrix}1&2\\0&-1\\4&1\end{psmallmatrix}$.
\item Verifique con matrices $2\times2$ concretas la identidad $(AB)^T=B^TA^T$.
\item Demuestre que $(A+B)^T=A^T+B^T$ para matrices de igual orden.
\end{listaejercicios}

\subsection{Inversas}
\begin{listaejercicios}
\item Mediante Gauss--Jordan, obtenga $A^{-1}$ para $A=\begin{psmallmatrix}3&1\\2&1\end{psmallmatrix}$ y compruebe el resultado.
\item Determine si $A=\begin{psmallmatrix}1&2\\3&6\end{psmallmatrix}$ tiene inversa; justifique mediante eliminación.
\item Calcule la inversa, si existe, de $A=\begin{psmallmatrix}1&0&0\\0&2&0\\0&0&-3\end{psmallmatrix}$.
\item Demuestre que $(A^{-1})^{-1}=A$ cuando $A$ es invertible.
\item Si $A$ y $B$ son invertibles, demuestre que $(AB)^{-1}=B^{-1}A^{-1}$.
\end{listaejercicios}

\subsection{Determinantes y cofactores}
\begin{listaejercicios}
\item Calcule $\begin{vmatrix}4&-2\\3&5\end{vmatrix}$.
\item Calcule $\begin{vmatrix}1&2&0\\0&3&1\\2&0&1\end{vmatrix}$ mediante Sarrus y expansión por cofactores.
\item Calcule el determinante de $\begin{psmallmatrix}2&1&0\\4&3&1\\0&2&3\end{psmallmatrix}$ mediante operaciones elementales.
\item Obtenga los menores y cofactores de la primera fila de $A=\begin{psmallmatrix}1&0&2\\3&1&-1\\0&2&4\end{psmallmatrix}$.
\item Sin expandir completamente, determine $\det\begin{psmallmatrix}1&2&3\\2&4&6\\0&1&1\end{psmallmatrix}$ y justifique.
\item Si $\det(A)=3$ para una matriz $3\times3$, determine $\det(2A)$ y $\det(A^T)$.
\item Calcule $\det(A)$ para $A=\begin{psmallmatrix}1&2&0&0\\0&3&1&0\\0&0&-2&5\\0&0&0&4\end{psmallmatrix}$.
\end{listaejercicios}

\subsection{Adjunta, rango y núcleo}
\begin{listaejercicios}
\item Mediante la fórmula de la adjunta, obtenga la inversa de $A=\begin{psmallmatrix}3&1\\2&1\end{psmallmatrix}$.
\item Determine $\adj(A)$ y $A^{-1}$, si existe, para $A=\begin{psmallmatrix}2&1&0\\1&2&1\\0&1&2\end{psmallmatrix}$.
\item Determine el rango y el núcleo de $A=\begin{psmallmatrix}1&2&3\\2&4&6\\0&1&1\end{psmallmatrix}$.
\item Obtenga el rango y el núcleo de $A=\begin{psmallmatrix}1&2&1\\0&1&1\end{psmallmatrix}$.
\item Compruebe que si $A$ es una matriz invertible de orden $n$, entonces $\Ker(A)=\{\vect0\}$.
\item Sea $A=\begin{psmallmatrix}1&1\\1&1\end{psmallmatrix}$ y $\vect b=\begin{psmallmatrix}2\\2\end{psmallmatrix}$. Describa todas las soluciones de $A\vect{x}=\vect b$ e indique por qué no puede usarse $A^{-1}\vect b$.
\end{listaejercicios}

\subsection{Ejercicios integradores tipo examen}
\begin{listaejercicios}
\item Sean $A=\begin{psmallmatrix}1&2\\-1&0\end{psmallmatrix}$ y $B=\begin{psmallmatrix}2&1\\0&-1\end{psmallmatrix}$. Calcule $AB$, $BA$, $\det(A)$ y determine si $A$ es invertible.
\item Resuelva $A\vect x=\vect b$ para $A=\begin{psmallmatrix}2&1\\1&1\end{psmallmatrix}$, $\vect b=\begin{psmallmatrix}7\\4\end{psmallmatrix}$ por (a) eliminación y (b) matriz inversa.
\item Sea $A=\begin{psmallmatrix}1&a\\2&4\end{psmallmatrix}$. Determine los valores de $a$ para los que $A$ no es invertible.
\item Determine $\rang(A)$ y $\Ker(A)$ para $A=\begin{psmallmatrix}1&1&0\\0&1&1\\1&2&1\end{psmallmatrix}$.
\item Calcule el determinante $\begin{vmatrix}1&0&2&0\\0&2&0&1\\0&0&-1&3\\0&0&0&4\end{vmatrix}$ y determine si la matriz correspondiente es invertible.
\end{listaejercicios}

%===============================================================
\appendix
\section{Soluciones seleccionadas paso a paso}
%===============================================================
Estas soluciones muestran el procedimiento, no sólo el resultado. Los números corresponden a la lista continua de ejercicios anterior.

\subsection*{Ejercicio E3: igualdad de matrices}
De las entradas respectivas:
\[
x-2=5,\qquad 3y=-6,\qquad z+1=4.
\]
Por tanto, $\boxed{x=7,\ y=-2,\ z=3}$.

\subsection*{Ejercicio E7: productos en distinto orden}
Para
\[
A=\begin{pmatrix}1&2\\0&1\end{pmatrix},\qquad
B=\begin{pmatrix}2&0\\1&3\end{pmatrix},
\]
calculamos
\[
AB=\begin{pmatrix}1(2)+2(1)&1(0)+2(3)\\0(2)+1(1)&0(0)+1(3)\end{pmatrix}
=\boxed{\begin{pmatrix}4&6\\1&3\end{pmatrix}},
\]
mientras que
\[
BA=\begin{pmatrix}2(1)+0(0)&2(2)+0(1)\\1(1)+3(0)&1(2)+3(1)\end{pmatrix}
=\boxed{\begin{pmatrix}2&4\\1&5\end{pmatrix}}.
\]
Así, $AB\ne BA$.

\subsection*{Ejercicio E12: inversa por Gauss--Jordan}
Para $A=\begin{psmallmatrix}3&1\\2&1\end{psmallmatrix}$ partimos de
\[
\left(\begin{array}{cc|cc}3&1&1&0\\2&1&0&1\end{array}\right)
\xrightarrow{R_1\leftarrow R_1-R_2}
\left(\begin{array}{cc|cc}1&0&1&-1\\2&1&0&1\end{array}\right)
\xrightarrow{R_2\leftarrow R_2-2R_1}
\left(\begin{array}{cc|cc}1&0&1&-1\\0&1&-2&3\end{array}\right).
\]
De aquí,
\[
\boxed{A^{-1}=\begin{pmatrix}1&-1\\-2&3\end{pmatrix}}.
\]
La verificación es
\[
\begin{pmatrix}3&1\\2&1\end{pmatrix}
\begin{pmatrix}1&-1\\-2&3\end{pmatrix}
=\begin{pmatrix}1&0\\0&1\end{pmatrix}.
\]

\subsection*{Ejercicio E17: determinante $2\times2$}
\[
\begin{vmatrix}4&-2\\3&5\end{vmatrix}=4(5)-(-2)(3)=20+6=\boxed{26}.
\]

\subsection*{Ejercicio E22: determinante de $2A$}
Para una matriz de orden $3$, multiplicar \emph{todas} sus entradas por $2$ equivale a multiplicar cada una de las tres filas por $2$. Luego
\[
\det(2A)=2^3\det(A)=8(3)=\boxed{24},\qquad
\det(A^T)=\det(A)=\boxed{3}.
\]

\subsection*{Ejercicio E24: inversa mediante adjunta}
Sea $A=\begin{psmallmatrix}3&1\\2&1\end{psmallmatrix}$. Su determinante es
\[
\det(A)=3(1)-1(2)=1.
\]
La matriz de cofactores es
\[
C=\begin{pmatrix}1&-2\\-1&3\end{pmatrix}.
\]
La adjunta es la transpuesta:
\[
\adj(A)=\begin{pmatrix}1&-1\\-2&3\end{pmatrix}.
\]
Finalmente,
\[
\boxed{A^{-1}=\frac{1}{\det(A)}\adj(A)
=\begin{pmatrix}1&-1\\-2&3\end{pmatrix}.}
\]
El resultado coincide con el ejercicio E12.

\subsection*{Ejercicio E27: rango y núcleo de una matriz rectangular}
Sea $A=\begin{psmallmatrix}1&2&1\\0&1&1\end{psmallmatrix}$. Eliminamos la entrada superior de la segunda columna:
\[
\begin{pmatrix}1&2&1\\0&1&1\end{pmatrix}
\xrightarrow{R_1\leftarrow R_1-2R_2}
\begin{pmatrix}1&0&-1\\0&1&1\end{pmatrix}.
\]
Hay dos pivotes: $\boxed{\rang(A)=2}$. El sistema homogéneo asociado es
\[
x_1-x_3=0,\qquad x_2+x_3=0.
\]
Tomando $x_3=t$:
\[
\boxed{\Ker(A)=\left\{t\begin{pmatrix}1\\-1\\1\end{pmatrix}:t\in\R\right\}.}
\]

\subsection*{Ejercicio E32: matriz con parámetro}
Sea
\[
A=\begin{pmatrix}1&a\\2&4\end{pmatrix}.
\]
La matriz no es invertible exactamente cuando
\[
\det(A)=4-2a=0.
\]
Despejando obtenemos $\boxed{a=2}$. Para cualquier $a\ne2$, la matriz sí es invertible.

\subsection*{Ejercicio E33: núcleo de una matriz $3\times3$}
Dada
\[
A=\begin{pmatrix}1&1&0\\0&1&1\\1&2&1\end{pmatrix},
\]
restamos la primera fila a la tercera:
\[
\begin{pmatrix}1&1&0\\0&1&1\\0&1&1\end{pmatrix}
\xrightarrow{R_3\leftarrow R_3-R_2}
\begin{pmatrix}1&1&0\\0&1&1\\0&0&0\end{pmatrix}.
\]
Así $\rang(A)=2$. El sistema equivalente es
\[
x_1+x_2=0,\qquad x_2+x_3=0.
\]
Si $x_3=t$, entonces $x_2=-t$ y $x_1=t$, con lo que
\[
\boxed{\Ker(A)=\left\{t\begin{pmatrix}1\\-1\\1\end{pmatrix}:t\in\R\right\}.}
\]

%===============================================================
\section{Síntesis para el repaso}
%===============================================================
\begin{center}
\renewcommand{\arraystretch}{1.35}
\begin{tabular}{>{\raggedright\arraybackslash}p{0.32\linewidth}>{\raggedright\arraybackslash}p{0.57\linewidth}}
\toprule
\textbf{Concepto} & \textbf{Regla o idea clave}\\\midrule
Suma de matrices & Mismo orden; entrada por entrada.\\
Producto $AB$ & Columnas de $A$ = filas de $B$; generalmente $AB\ne BA$.\\
Transpuesta & $(AB)^T=B^TA^T$.\\
Inversa & $AA^{-1}=A^{-1}A=I$; sólo para matrices cuadradas invertibles.\\
Determinante $2\times2$ & $\det\begin{psmallmatrix}a&b\\c&d\end{psmallmatrix}=ad-bc$.\\
Cofactor & $C_{ij}=(-1)^{i+j}M_{ij}$.\\
Adjunta & $\adj(A)=C^T$.\\
Inversa por adjunta & $A^{-1}=\adj(A)/\det(A)$, si $\det(A)\ne0$.\\
Rango & Número de pivotes.\\
Núcleo & Soluciones de $A\vect x=\vect0$.\\\bottomrule
\end{tabular}
\end{center}

\section{Bibliografía de referencia}
\begin{itemize}
\item Grossman, S. I. (2012). \emph{Álgebra Lineal}. 7a edición. McGraw-Hill.
\item Lay, D. C. (2013). \emph{Álgebra lineal para cursos con enfoque por competencias}. Pearson.
\item Kolman, B. (2013). \emph{Álgebra Lineal}. Pearson Educación.
\item Poole, D. (2011). \emph{Álgebra lineal: una introducción moderna}. 3a edición. Cengage Learning.
\end{itemize}

\vspace{0.8em}\begin{center}\textit{Fin de la primera propuesta de la Unidad 3.}\end{center}
\end{document}
